%%================================================
%% Appendix E — UX instrument summary
%%================================================
\chapter{ภาคผนวก จ: แบบสอบถามและขั้นตอนการประเมิน UX ของผู้ดูแล}
\label{app:ux_instrument}

ภาคผนวกนี้สรุปโครงสร้างเครื่องมือประเมินประสบการณ์ผู้ใช้ฝั่งผู้ดูแลที่ใช้ในบทที่~\ref{chapter4} โดยไม่เปิดเผยข้อมูลระบุตัวบุคคลของผู้ตอบหรือข้อความดิบทั้งชุดจากช่องเปิด

\section*{การเชิญชวนและลักษณะกลุ่มตัวอย่าง}

ผู้เข้าร่วมได้รับเชิญผ่านช่องทางภายในโครงการ เช่น ผู้ดูแลหรือผู้เกี่ยวข้องที่เข้าถึงระบบต้นแบบในพื้นที่ทดสอบและช่วงเวลาที่กำหนด การตอบแบบสอบถามเป็นโดยสมัครใจ ไม่มีการบังคับ และแต่ละลิงก์ถูกตั้งให้ตอบได้ครั้งเดียวตามการตั้งค่าของ Google Form ผลที่อ้างในเล่มจำกัดอยู่ที่ผู้ตอบที่กรอกครบในช่วงหน้าต่างการประเมินที่ระบุในบทที่~\ref{chapter4} จึงเป็นกลุ่มตัวอย่างเชิงความสะดวก ไม่ได้สุ่มจากประชากรผู้ดูแลโดยรวม และไม่ควรนำไปอนุมานในระดับประชากร

\section*{แหล่งที่มาและขอบเขตการใช้งาน}

แบบสอบถามจัดทำเป็น Google Form ภาษาไทย และข้อมูลที่ใช้ในบทที่~\ref{chapter4} ถูกรวบรวมจากไฟล์ผลตอบกลับเดียวกันที่ใช้สรุปตารางคะแนน Likert ช่องทางใช้งาน และธีมเชิงคุณภาพ แบบสอบถามนี้ออกแบบเพื่อประเมินต้นแบบในบริบทการใช้งานจริงย่อส่วน ไม่ได้ออกแบบเป็นเครื่องมือมาตรฐานเพื่ออนุมานผลในระดับประชากรหรือเปรียบเทียบกับเครื่องมือทางจิตวิทยามาตรฐานโดยตรง

\section*{จริยธรรมและความเป็นส่วนตัว}

แบบสอบถามไม่เผยแพร่ข้อมูลที่ระบุตัวบุคคลโดยตรงในเล่ม เช่น ชื่อผู้ตอบหรือข้อมูลส่วนบุคคลอื่น ๆ ความเห็นปลายเปิดถูกสรุปเป็นกลุ่มธีมในบทที่~\ref{chapter4} เพื่อลดความเสี่ยงจากการโยงย้อนข้อความกลับไปยังผู้ตอบรายบุคคล และการใช้ Google Form อยู่ภายใต้ข้อตกลงของผู้ให้บริการแบบฟอร์มตามสภาพแวดล้อมการเก็บข้อมูลของโครงการ

\section*{มิติการให้คะแนนแบบ Likert}

แบบสอบถามใช้มาตราส่วน \(1\)--\(5\) โดย \(5\) หมายถึงความเห็นเชิงบวกหรือความพึงพอใจในระดับสูงสุดตามข้อความของแต่ละข้อ มิติที่ใช้ในเล่มมี 5 ข้อหลัก และเป็นฐานของตารางที่~\ref{tab:ch4_ux_summary}

\begin{enumerate}
    \item การใช้งานหน้าจอผู้ดูแลโดยรวมมีความง่ายต่อการใช้
    \item ปุ่ม เมนู หรือจุดที่ต้องกดไม่มากหรือซับซ้อนเกินไป
    \item หากไม่มีผู้สอน สามารถลองใช้และเข้าใจเบื้องต้นด้วยตนเองได้
    \item คำอธิบายและป้ายบนหน้าจออ่านแล้วเข้าใจได้ง่าย
    \item ระบบช่วยการทำงานในฐานะผู้ดูแลได้ในระดับที่พึงพอใจ
\end{enumerate}

\noindent\textbf{หมายเหตุ:} ข้อที่สองเป็นคำถามเชิงความซับซ้อนแต่ยังใช้สเกลเดียวกับข้ออื่น การตีความเชิงปริมาณจึงควรอ่านควบคู่กับข้อความอธิบายและผลเชิงคุณภาพในบทที่~\ref{chapter4}

\section*{ฟิลด์บทบาท ช่องทางใช้งาน และข้อความปลายเปิด}

นอกจากข้อ Likert แบบสอบถามยังเก็บข้อมูลบทบาทของผู้ตอบ ช่องทางที่คาดว่าจะใช้งานบ่อยที่สุด เช่น เว็บหรือมือถือ และคำถามปลายเปิดเกี่ยวกับฟีเจอร์ที่มีประโยชน์ ความยากในการใช้งาน และข้อเสนอแนะในการปรับปรุง ข้อมูลเหล่านี้เป็นฐานของตารางที่~\ref{tab:ch4_ux_channel_role} และ~\ref{tab:ch4_ux_themes}

\section*{ความสำคัญต่อการตีความผล}

การอ่านผล UX ในบทที่~\ref{chapter4} ควรคำนึงว่าแบบสอบถามชุดนี้มุ่งใช้เพื่อสะท้อนทิศทางการออกแบบของต้นแบบ เช่น ความจำเป็นของป้ายภาษาไทย การแจ้งเตือนที่ชัดเจน และความสำคัญของช่องทางมือถือ มากกว่าการสรุปคุณภาพการใช้งานในเชิงประชากรศาสตร์ ภาคผนวกนี้จึงช่วยให้ผู้อ่านเข้าใจกรอบการตีความข้อมูล UX ได้ชัดขึ้นและเชื่อมโยงข้อค้นพบกับบริบทการเก็บข้อมูลจริง
